\section{How to implement these ‘Coding Conventions’ in your Project?}\label{sec:HowToImplement}
First any coding conventions are better than none! So if your team is already adhering to some coding conventions, the first hurdle had been taken already.

Adjust this document that it will incorporate the already accepted guidelines and recommendations; this may require that you remove conflicting recommendations and guidelines – don't care, just do it!

The biggest blocker in this case could be that too much code needs to be modified to match with the coding conventions, meaning tedious work no one likes to do. In this regard, you should also have a look to the chapter \tqfullvref{sec:CodingForSustaining}.

But the most simple case is a new team starting over with a new project from scratch. Just present the document as is~…

Experienced programmers know about the benefits of having detailed guard\-rails they can follow – in particular when they know that the other team members will also follow the same guardrails.

Nevertheless, you may have one or the other hacker in your team that complains that this kind of guardrails are suffocating their creativity.

Perhaps this helps: I was always interested in fotography, and therefore I attended some classes about the topic. There the question was raised what are bad pictures, good pictures and exceptionally artistic pictures, and the professor explained it like this: Fotography has a whole bunch of rules on how to compose a good picture, and you do not even need to know one of these rules to distinguish a bad shot from a good fotography: the bad ones will break the rules.

But the really artistic pictures always break \textit{one} or \textit{the other} rule, but \textit{with intention}. The fotographer's creativity is to choose the rule to break and how – but this first means that they have to know each of the rules and how to follow them (usually referred to as ‘knowing the craft’), and second, that there is a visible intention.

But in most cases the client just wants to get only good pictures.

Same in software development: only few clients pay for a piece of art; instead they expect good craftsmanship.

If this explanation does not help, and you are in the proper position: replace the hacker.

%==============================================================================
% End of file!!
%==============================================================================
